\section*{Hoofdstuk \ref{ch:b1:fractions} --- Rationale breuken}

\begin{solution}{exo:b1:fractions:1}
$\dfrac{1}{X^2 - 1}$: enkelvoudige \omterm{def:b1:fractions:field}{polen} $\pm 1$; afdekken:
$\dfrac{1/2}{X-1} - \dfrac{1/2}{X+1}$.

$\dfrac{X}{X^2 - 3X + 2} = \dfrac{X}{(X-1)(X-2)}$: afdekken geeft
$\frac{1}{1-2} = -1$ in $1$ en $\frac{2}{2-1} = 2$ in $2$:
$\dfrac{-1}{X-1} + \dfrac{2}{X-2}$.

$\dfrac{X^2+1}{X(X-1)}$: de graad is $0$, dus er is een geheel
deel: door te delen, $X^2 + 1 = (X^2 - X) + (X + 1)$, dus $F = 1 +
\frac{X+1}{X(X-1)}$. Afdekken op de rest: $\frac{1}{-1} = -1$
in $0$, $\frac{2}{1} = 2$ in $1$:
\[
F = 1 - \frac{1}{X} + \frac{2}{X-1} .
\]
\end{solution}

\begin{solution}{exo:b1:fractions:2}
$\dfrac{1}{X(X^2+1)}$: vorm $\frac aX + \frac{bX + c}{X^2 + 1}$.
Afdekken in $0$: $a = 1$. Limiet van $xF$: $0 = a + b$, dus $b = -1$.
Evaluatie in $X = 1$: $\frac12 = 1 + \frac{c - 1}{2}$, dus $c = 0$:
\[
\frac{1}{X(X^2+1)} = \frac 1X - \frac{X}{X^2+1} .
\]

$\dfrac{X^3}{X^2+X+1}$: deling: $X^3 = (X^2+X+1)(X - 1) + 1$, dus
\[
\frac{X^3}{X^2+X+1} = X - 1 + \frac{1}{X^2 + X + 1} ,
\]
reeds in reële gesplitste vorm (de kwadratische \omterm{def:b1:poly:def}{veelterm} heeft
negatieve discriminant).
\end{solution}

\begin{solution}{exo:b1:fractions:3}
$\dfrac{1}{X^2(X-1)}$: vorm $\frac{a}{X^2} + \frac bX +
\frac{c}{X-1}$. Afdekken in de dubbele \omterm{def:b1:fractions:field}{pool} $0$:
$a = \bigl[\frac{1}{X-1}\bigr]_{0} = -1$. Afdekken in $1$: $c = 1$.
Limiet van $xF$: $0 = b + c$, dus $b = -1$:
\[
\frac{1}{X^2(X-1)} = -\frac{1}{X^2} - \frac1X + \frac{1}{X-1} .
\]

$\dfrac{X+1}{(X-1)^3}$: met $Y = X - 1$ is de teller $Y + 2$:
\[
\frac{Y + 2}{Y^3} = \frac{1}{Y^2} + \frac{2}{Y^3}
= \frac{1}{(X-1)^2} + \frac{2}{(X-1)^3} .
\]
\end{solution}

\begin{solution}{exo:b1:fractions:4}
De \omterm{def:b1:fractions:field}{polen} zijn de vierde \omterm{def:b1:complex:unity}{eenheidswortels} $1, \iu, -1, -\iu$, alle
enkelvoudig. Formule voor een enkelvoudige \omterm{def:b1:fractions:field}{pool} met $B' = 4X^3$: de
coëfficiënt in $a$ is $\frac{1}{4a^3} = \frac{a}{4a^4} = \frac a4$
(gebruikmakend van $a^4 = 1$). Dus, over $\C$:
\[
\frac{1}{X^4 - 1}
= \frac{1/4}{X - 1} - \frac{1/4}{X + 1}
+ \frac{\iu/4}{X - \iu} - \frac{\iu/4}{X + \iu} .
\]
Het toegevoegd complexe paar groeperen (gemeenschappelijke noemer
$X^2 + 1$): $\frac{\iu}{4}\bigl(\frac{1}{X-\iu} - \frac{1}{X+\iu}\bigr) =
\frac{\iu}{4}\cdot\frac{2\iu}{X^2+1} = \frac{-1/2}{X^2+1}$. Over
$\R$:
\[
\frac{1}{X^4 - 1}
= \frac{1/4}{X-1} - \frac{1/4}{X+1} - \frac{1/2}{X^2 + 1} .
\]
\emph{Controle in $X = 0$:} $-1 = -\frac14 - \frac14 - \frac12$.
\end{solution}

\begin{solution}{exo:b1:fractions:5}
$\dfrac{X^2}{(X^2+1)^2} = \dfrac{(X^2 + 1) - 1}{(X^2+1)^2} =
\dfrac{1}{X^2+1} - \dfrac{1}{(X^2+1)^2}$.

Bijgevolg een primitieve:
\[
\int \frac{x^2\,\dd x}{(x^2+1)^2}
= \arctan x - \frac12\Bigl(\arctan x + \frac{x}{x^2+1}\Bigr) + C
= \frac12\arctan x - \frac{x}{2(x^2+1)} + C .
\]
\end{solution}

\begin{solution}{exo:b1:fractions:6}
De \omterm{def:b1:fractions:field}{polen} $0, -1, \dots, -n$ zijn enkelvoudig. Afdekken in $-k$:
\[
c_k = \frac{n!}{\prod_{j \neq k} (-k + j)}
= \frac{n!}{\bigl(\prod_{j=0}^{k-1}(j - k)\bigr)
\bigl(\prod_{j=k+1}^{n}(j-k)\bigr)}
= \frac{n!}{(-1)^k k!\,(n-k)!} = (-1)^k \binom nk .
\]
Dus
\[
\frac{n!}{X(X+1)\cdots(X+n)}
= \sum_{k=0}^{n} \frac{(-1)^k \binom nk}{X + k} .
\]
(Plausibiliteitscontrole voor $n = 1$: $\frac{1}{X(X+1)} = \frac1X -
\frac{1}{X+1}$.)
\end{solution}

\begin{solution}{exo:b1:fractions:7}
$P = X^n - 1$ heeft de $n$ enkelvoudige wortels $\omega^k$
(\cref{thm:b1:complex:roots}), dus de identiteit van de logaritmische
afgeleide van \cref{exo:b1:poly:11} luidt
\[
\sum_{k=0}^{n-1} \frac{1}{X - \omega^k}
= \frac{P'(X)}{P(X)} = \frac{n X^{n-1}}{X^n - 1} .
\]
In $X = 2$, $n = 4$: rechterlid $= \frac{4 \times 8}{15} =
\frac{32}{15}$. Linkerlid: $\frac{1}{2-1} + \frac{1}{2+1} +
\frac{1}{2 - \iu} + \frac{1}{2 + \iu} = 1 + \frac13 + \frac{4}{5} =
\frac{15 + 5 + 12}{15} = \frac{32}{15}$, gebruikmakend van
$\frac{1}{2-\iu} + \frac{1}{2+\iu} = \frac{4}{5}$.
\end{solution}

\begin{solution}{exo:b1:fractions:8}
Afdekken: $\dfrac{1}{k(k+1)(k+2)} = \dfrac{1/2}{k} - \dfrac{1}{k+1} +
\dfrac{1/2}{k+2}$. Herschrijf als een telescoperend verschil:
\[
\frac{1}{k(k+1)(k+2)}
= \frac12\Bigl(\frac{1}{k(k+1)} - \frac{1}{(k+1)(k+2)}\Bigr),
\]
(ontwikkel ter controle --- of trek de twee splitsingen van elkaar
af). Sommeren:
\[
S_n = \frac12\Bigl(\frac{1}{1 \times 2} - \frac{1}{(n+1)(n+2)}\Bigr)
= \frac14 - \frac{1}{2(n+1)(n+2)}
\xrightarrow[n \to \infty]{} \frac14 .
\]
\end{solution}

\begin{solution}{exo:b1:fractions:9}
Splits, voor $0 \leq m \leq n - 1$, de breuk $\frac{X^m}{P}$
(graad $m - n \leq -1$, enkelvoudige \omterm{def:b1:fractions:field}{polen}): de formule voor een
enkelvoudige \omterm{def:b1:fractions:field}{pool} geeft
\[
\frac{X^m}{P} = \sum_{i=1}^{n} \frac{x_i^m / P'(x_i)}{X - x_i} .
\]
Vermenigvuldig met $X$ en laat $X \to +\infty$: het linkerlid neigt
naar de limiet van $X^{m+1}/P$, die $0$ is als $m \leq n - 2$ en $1$
als $m = n - 1$ ($P$ is \omterm{def:b1:poly:def}{monisch} van graad $n$); het rechterlid neigt
naar $\sum_i \frac{x_i^m}{P'(x_i)}$. Bijgevolg
\[
\sum_{i=1}^{n} \frac{x_i^{m}}{P'(x_i)} =
\begin{cases}
0 & \text{voor } 0 \leq m \leq n-2,\\
1 & \text{voor } m = n - 1,
\end{cases}
\]
wat beide aangekondigde identiteiten bevat ($m = 0$ vereist $n \geq
2$). (Interpretatie via \cref{thm:b1:poly:lagrange}: deze sommen zijn
de leidende coëfficiënten van de Lagrange-interpolanten van $X^m$, en
een \omterm{def:b1:poly:def}{veelterm} van graad $\leq n-1$ in $n$ punten interpoleren
reproduceert hem exact.)
\end{solution}

\begin{solution}{exo:b1:fractions:10}
Vorm $\frac aX + \frac b{X+1} + \frac c{(X+1)^2}$. Afdekken in
$0$: $a = 1$; afdekken in de dubbele \omterm{def:b1:fractions:field}{pool}: $c = \bigl[\frac1X
\bigr]_{X=-1} = -1$; limiet van $xF(x)$ in oneindig: $0 = a + b$, dus
$b = -1$:
\[
\frac{1}{X(X+1)^2} = \frac1X - \frac1{X+1} - \frac1{(X+1)^2} .
\]
Sommeren voor $n = 1, \dots, N$: de eerste twee bouwstenen telescoperen
tot $1 - \frac1{N+1}$, en de derde levert $-\sum_{k=2}^{N+1}
\frac1{k^2}$. Als $N \to \infty$ en gebruikmakend van $\sum_{k\geq1}
\frac1{k^2} = \frac{\pi^2}6$:
\[
\sum_{n=1}^{\infty}\frac{1}{n(n+1)^2}
= 1 - \Bigl(\frac{\pi^2}6 - 1\Bigr) = 2 - \frac{\pi^2}6
\approx 0.355 .
\]
\end{solution}

\begin{solution}{exo:b1:fractions:11}
In de variabele $Y$: $\frac1{(Y+1)(Y+4)} = \frac{1/3}{Y+1} -
\frac{1/3}{Y+4}$ (afdekken in $-1$ en $-4$), dus
\[
\frac{1}{(X^2+1)(X^2+4)}
= \frac13\,\frac1{X^2+1} - \frac13\,\frac1{X^2+4} .
\]
Evenzo $\frac{Y}{(Y+1)(Y+4)} = \frac{-1/3}{Y+1} +
\frac{4/3}{Y+4}$, dus
\[
\frac{X^2}{(X^2+1)(X^2+4)}
= -\frac13\,\frac1{X^2+1} + \frac43\,\frac1{X^2+4} .
\]
(Controle in $X = 0$: $0 = \frac13(-1 + 1)$.) Dit zijn reeds de reële
splitsingen: de tellers boven de onontbindbare kwadratische \omterm{def:b1:poly:def}{veeltermen}
blijken constanten te zijn.
\end{solution}

\begin{solution}{exo:b1:fractions:12}
\begin{enumerate}
  \item Op elk interval dat de \omterm{def:b1:fractions:field}{polen} vermijdt, $F'(x) = -\sum_i
        \frac{c_i}{(x - p_i)^2} < 0$: strikt dalend. Als $x
        \to \pm\infty$ neigt elke bouwsteen naar $0$: $F \to 0$, van
        boven in $+\infty$ (alle bouwstenen daar positief) en van
        onder in $-\infty$. Als $x \to p_i^+$ blaast de bouwsteen
        $\frac{c_i}{x - p_i}$ op naar $+\infty$ en de andere
        blijven begrensd: $F \to +\infty$; evenzo $F \to -\infty$ als
        $x \to p_i^-$.
  \item Fixeer $\lambda > 0$. Op $\intoo{-\infty}{p_1}$: $F$
        daalt van $0^-$ naar $-\infty$, dus $F < 0 < \lambda$:
        geen oplossing. Op elk $\intoo{p_i}{p_{i+1}}$ ($1 \leq i
        \leq r-1$): $F$ daalt van $+\infty$ naar $-\infty$,
        neemt dus de waarde $\lambda$ precies één keer aan
        (tussenwaardeeigenschap plus strikte monotonie). Op
        $\intoo{p_r}{+\infty}$: $F$ daalt van $+\infty$ naar
        $0^+$, opnieuw precies één oplossing. Totaal: precies $r$
        oplossingen, verstrengeld met de \omterm{def:b1:fractions:field}{polen}.
\end{enumerate}
\end{solution}

\begin{solution}{pb:b1:fractions:1}
\textbf{1.} Afdekken: $\frac1{X(X+1)} = \frac1X - \frac1{X+1}$.
De som telescopeert:
\[
\sum_{n=1}^{N}\Bigl(\frac1n - \frac1{n+1}\Bigr)
= 1 - \frac1{N+1} \longrightarrow 1 .
\]

\textbf{2.} $\frac1{X(X+2)} = \frac{1/2}X - \frac{1/2}{X+2}$.
Bij het sommeren overleven de termen $\frac1n$ voor $n = 1, 2$ en de
termen $-\frac1{n+2}$ voor $n = N-1, N$:
\[
\sum_{n=1}^{N}\frac1{n(n+2)}
= \frac12\Bigl(1 + \frac12 - \frac1{N+1} - \frac1{N+2}\Bigr)
\longrightarrow \frac34 .
\]

\textbf{3.} Als $F(X) = G(X) - G(X+1)$, dan $\sum_{n=1}^N F(n) =
\sum_{n=1}^N\bigl(G(n) - G(n+1)\bigr) = G(1) - G(N+1)$: alle
tussenliggende waarden heffen paarsgewijs op. Voor $F =
\frac1{X(X+1)(X+2)}$ is de getuige $G(X) = \frac1{2X(X+1)}$:
\[
G(X) - G(X+1)
= \frac{(X+2) - X}{2X(X+1)(X+2)} = F(X) ,
\]
dus $\sum_{n=1}^N F(n) = \frac14 - \frac1{2(N+1)(N+2)} \to
\frac14$, de waarde van \cref{exo:b1:fractions:8}.

\textbf{4.} Breng het rechterlid op de gemeenschappelijke noemer
$X(X+1)\cdots(X+k)$:
\[
\frac1k\cdot\frac{(X + k) - X}{X(X+1)\cdots(X+k)}
= \frac{1}{X(X+1)\cdots(X+k)} ,
\]
wat de identiteit is. Dus $F_k(X) = \frac1k\bigl(G_k(X) -
G_k(X+1)\bigr)$ met $G_k(X) = \frac1{X(X+1)\cdots(X+k-1)}$, en het
mechanisme van vraag 3 geeft
\[
\sum_{n=1}^{N}\frac1{n(n+1)\cdots(n+k)}
= \frac1k\Bigl(\frac1{k!} - G_k(N+1)\Bigr)
\longrightarrow \frac1{k\cdot k!} ,
\]
aangezien $G_k(1) = \frac1{k!}$ en $G_k(N+1) \to 0$. Voor $k = 2$:
$\frac1{2\cdot2} = \frac14$, in overeenstemming met vraag 3.

\textbf{5.} \cref{exo:b1:fractions:6} geeft
$\frac{n!}{X(X+1)\cdots(X+n)} = \sum_{j=0}^n
\frac{(-1)^j\binom nj}{X + j}$. Evalueer in $X = 1$: het linkerlid
is $\frac{n!}{(n+1)!} = \frac1{n+1}$, het rechterlid is
$\sum_j(-1)^j\binom nj\frac1{1+j}$: eerste identiteit. In $X = 2$:
het linkerlid is $\frac{n!}{2\cdot3\cdots(n+2)} =
\frac{n!\cdot1}{(n+2)!} = \frac1{(n+1)(n+2)}$, het rechterlid
$\sum_j(-1)^j\binom nj\frac1{2+j}$: tweede identiteit.

\textbf{6.} De \omterm{def:b1:fractions:field}{polen} $\omega^k$ zijn enkelvoudig, en de formule voor
een enkelvoudige \omterm{def:b1:fractions:field}{pool} van \cref{met:b1:fractions:compute} geeft de
coëfficiënt
\[
\frac{1}{\bigl(nX^{n-1}\bigr)_{X = \omega^k}}
= \frac{1}{n\,\omega^{k(n-1)}}
= \frac{\omega^k}{n\,\omega^{kn}} = \frac{\omega^k}n ,
\]
gebruikmakend van $\omega^{kn} = 1$. Bijgevolg $\frac1{X^n-1} =
\frac1n\sum_k \frac{\omega^k}{X - \omega^k}$.

\textbf{7.} Voor $n = 2$ ($\omega = -1$): $\frac12\bigl(
\frac1{X-1} - \frac1{X+1}\bigr) = \frac12\cdot\frac{2}{X^2-1} =
\frac1{X^2-1}$: correct. De coëfficiënten sommeren tot
$\frac1n\sum_k\omega^k = 0$ voor $n \geq 2$
(\cref{prop:b1:complex:sumroots}). Ze moeten wel: $x\,F(x) \to
\sum_k c_k$ als $x \to \infty$ voor elke splitsing met enkelvoudige
\omterm{def:b1:fractions:field}{polen}, terwijl hier $xF(x) = \frac{x}{x^n-1} \to 0$ aangezien $n \geq
2$.

\textbf{8.} Afdekken voor $\frac{X^{n-1}}{X^n - 1}$ in $\omega^k$:
$\frac{A(\omega^k)}{B'(\omega^k)} =
\frac{\omega^{k(n-1)}}{n\omega^{k(n-1)}} = \frac1n$, dus
$\frac{X^{n-1}}{X^n-1} = \frac1n\sum_k\frac1{X - \omega^k}$ ---
opnieuw de identiteit van \cref{exo:b1:fractions:7}.

\textbf{9.} Met $c = \omega^k$, $\conj c = \omega^{n-k}$ en
$c\conj c = 1$, $c + \conj c = 2\cos\theta_k$:
\[
\frac{c}{X - c} + \frac{\conj c}{X - \conj c}
= \frac{c(X - \conj c) + \conj c(X - c)}
{(X - c)(X - \conj c)}
= \frac{2\cos\theta_k\,X - 2}{X^2 - 2\cos\theta_k\,X + 1} .
\]
Door $k$ met $n - k$ te groeperen in vraag 6: voor oneven $n$,
\[
\frac1{X^n - 1} = \frac1n\Biggl(\frac1{X-1} +
\sum_{k=1}^{(n-1)/2}
\frac{2\cos\theta_k X - 2}{X^2 - 2\cos\theta_k X + 1}\Biggr) ;
\]
voor even $n$ draagt de extra zelf-gekoppelde \omterm{def:b1:fractions:field}{pool} $\omega^{n/2} = -1$
$\frac{-1}{X+1}$ bij binnen het haakje en loopt de somover de paren
tot $\frac n2 - 1$.

\textbf{10.} $n = 4$: $\theta_1 = \frac\pi2$, $\cos\theta_1 = 0$,
dus de paarterm is $\frac{-2}{X^2+1}$ en
\[
\frac1{X^4-1} = \frac14\Bigl(\frac1{X-1} - \frac1{X+1}
- \frac2{X^2+1}\Bigr) ,
\]
de splitsing van \cref{exo:b1:fractions:4}.

\textbf{11.} $P = \prod_{k=1}^{n-1}(X - \omega^k)$ (deel $X^n -
1$ door $X - 1$), dus door \cref{exo:b1:poly:11}, $\frac{P'}P =
\sum_{k=1}^{n-1}\frac1{X-\omega^k}$. In $X = 1$: $P(1) = n$ en
$P'(1) = \sum_{j=1}^{n-1}j = \frac{n(n-1)}2$, waaruit
\[
\sum_{k=1}^{n-1}\frac1{1 - \omega^k}
= \frac{P'(1)}{P(1)} = \frac{n-1}2 .
\]

\textbf{12.} Halvehoek: $1 - \eu^{\iu\theta} =
-2\iu\sin\frac\theta2\,\eu^{\iu\theta/2}$, dus
\[
\frac1{1 - \eu^{\iu\theta}}
= \frac{\eu^{-\iu\theta/2}}{-2\iu\sin\frac\theta2}
= \frac{\iu\cos\frac\theta2 + \sin\frac\theta2}
{2\sin\frac\theta2}
= \frac12 + \frac\iu2\,
\frac{\cos\frac\theta2}{\sin\frac\theta2} .
\]
Met $\theta = \theta_k = \frac{2k\pi}n$: $\frac1{1-\omega^k} =
\frac12 + \frac\iu2\cot\frac{k\pi}n$. Sommeren over $k = 1, \dots,
n-1$ en vergelijken met de reële waarde $\frac{n-1}2$ van
vraag 11: de reële delen verklaren reeds alles, dus
$\sum_k\cot\frac{k\pi}n = 0$ --- zoals de symmetrie $\cot\frac{(n-k)
\pi}n = -\cot\frac{k\pi}n$ eveneens aantoont.

\textbf{13.} Door $\frac{P'}P = \sum_k\frac1{X -
\omega^k}$ te differentiëren:
\[
\frac{P''}P - \Bigl(\frac{P'}P\Bigr)^2
= -\sum_{k=1}^{n-1}\frac1{(X - \omega^k)^2} .
\]
In $X = 1$: $P''(1) = \sum_{j=2}^{n-1}j(j-1) = 2\binom n3 =
\frac{n(n-1)(n-2)}3$ (hockeystick-identiteit, of inductie), dus
\[
\sum_{k=1}^{n-1}\frac1{(1 - \omega^k)^2}
= \Bigl(\frac{n-1}2\Bigr)^2 - \frac{(n-1)(n-2)}3
= \frac{(n-1)\bigl(3(n-1) - 4(n-2)\bigr)}{12}
= \frac{(n-1)(5-n)}{12} .
\]

\textbf{14.} Kwadrateren van de formule van vraag 12, met $c_k =
\cot\frac{k\pi}n$:
\[
\frac1{(1-\omega^k)^2}
= \Bigl(\frac12 + \frac\iu2 c_k\Bigr)^2
= \frac14 - \frac{c_k^2}4 + \frac\iu2\,c_k .
\]
Sommeren en gebruikmakend van $\sum c_k = 0$ (vraag 12) en vraag 13:
$\frac{n-1}4 - \frac14\sum_k c_k^2 = \frac{(n-1)(5-n)}{12}$, dus
\[
\sum_{k=1}^{n-1}\cot^2\frac{k\pi}n
= (n-1) - \frac{(n-1)(5-n)}3 = \frac{(n-1)(n-2)}3 .
\]
Dan geeft $\frac1{\sin^2t} = 1 + \cot^2t$
$\sum_k\frac1{\sin^2\frac{k\pi}n} = (n-1) +
\frac{(n-1)(n-2)}3 = \frac{n^2-1}3$.

\textbf{15.} $n = 3$: $\cot^2\frac\pi3 + \cot^2\frac{2\pi}3 =
\frac13 + \frac13 = \frac23 = \frac{2\cdot1}3$; en
$\frac1{\sin^2}$ sommeert tot $\frac43 + \frac43 = \frac83 =
\frac{9-1}3$. $n = 4$: $1 + 0 + 1 = 2 = \frac{3\cdot2}3$; en $2
+ 1 + 2 = 5 = \frac{16-1}3$. Beide formules kloppen.

\textbf{16.} Voor $t \in \intoo0{\frac\pi2}$ levert de klassieke
vergelijking $\sin t < t < \tan t$ (oppervlakte- of
convexiteitsargument, bekend van de middelbare school), door
inverses te nemen, $\cot t < \frac1t < \frac1{\sin t}$, alle drie
daar positief; kwadrateren behoudt de volgorde. Met $t =
\frac{k\pi}n$, $1 \leq k \leq m$, $n = 2m+1$ (zodat $t <
\frac\pi2$):
\[
\cot^2\frac{k\pi}n < \frac{n^2}{k^2\pi^2} <
\frac1{\sin^2\frac{k\pi}n} .
\]

\textbf{17.} Door de symmetrieën $\cot^2\frac{(n-k)\pi}n =
\cot^2\frac{k\pi}n$ en evenzo voor $\sin^2$, halveren de sommen van
vraag 14: $\sum_{k=1}^{m}\cot^2\frac{k\pi}n = \frac{(n-1)(n-2)}6 =
\frac{m(2m-1)}3$ en $\sum_{k=1}^m\frac1{\sin^2\frac{k\pi}n} =
\frac{n^2-1}6 = \frac{2m(m+1)}3$. Sommeren van vraag 16 over $k =
1, \dots, m$ en vermenigvuldigen met $\frac{\pi^2}{n^2}$:
\[
\frac{\pi^2}{(2m+1)^2}\cdot\frac{m(2m-1)}3
\;<\; \sum_{k=1}^{m}\frac1{k^2} \;<\;
\frac{\pi^2}{(2m+1)^2}\cdot\frac{2m(m+1)}3 .
\]
Beide grenzen neigen naar $\frac{\pi^2}6$ als $m \to \infty$ (de
verhoudingen $\frac{m(2m-1)}{(2m+1)^2}$ en $\frac{2m(m+1)}{(2m+1)^2}$
neigen beide naar $\frac12$), dus door het inknellen convergeren de
stijgende deelsommen en
\[
\sum_{k=1}^{\infty}\frac1{k^2} = \frac{\pi^2}6 .
\]

\textbf{18.} Kwadrateer $\frac1{X^2-1} = \frac12\bigl(\frac1{X-1} -
\frac1{X+1}\bigr)$:
\[
\frac1{(X^2-1)^2} = \frac14\Bigl(\frac1{(X-1)^2} +
\frac1{(X+1)^2}\Bigr) - \frac12\cdot\frac1{(X-1)(X+1)} ,
\]
en hersplits de kruisterm $\frac1{(X-1)(X+1)} =
\frac12\bigl(\frac1{X-1} - \frac1{X+1}\bigr)$ om de vermelde
vorm te bekomen. In $X = 0$: linkerlid $1$; rechterlid $\frac14(1 + 1) -
\frac14(-1 - 1) = \frac12 + \frac12 = 1$.

\textbf{19.} Sommeer vraag 18 over $n \geq 2$. Met $S =
\sum_{k\geq1}\frac1{k^2} = \frac{\pi^2}6$: $\sum_{n\geq2}
\frac1{(n-1)^2} = S$; $\sum_{n\geq2}\frac1{(n+1)^2} = S - 1 -
\frac14$; en de telescoop met tussenruimte twee $\sum_{n\geq2}\bigl(
\frac1{n-1} - \frac1{n+1}\bigr) = 1 + \frac12 = \frac32$. Bijgevolg
\[
\sum_{n=2}^{\infty}\frac1{(n^2-1)^2}
= \frac14\Bigl(2S - \frac54\Bigr) - \frac14\cdot\frac32
= \frac S2 - \frac{11}{16}
= \frac{\pi^2}{12} - \frac{11}{16} \approx 0.135 .
\]
Numeriek: $\frac19 + \frac1{64} + \frac1{225} + \frac1{576} +
\frac1{1225} + \dots \approx 0.1111 + 0.0156 + 0.0044 + 0.0017 +
0.0008 + \dots \approx 0.135$: consistent.

\textbf{20.} Differentieer de identiteit van vraag 8
$\sum_k\frac1{X-\omega^k} = \frac{nX^{n-1}}{X^n-1}$:
\[
\sum_{k=0}^{n-1}\frac1{(X - \omega^k)^2}
= \frac{n^2X^{2n-2} - n(n-1)X^{n-2}(X^n - 1)}{(X^n - 1)^2} .
\]
Controle in $n = 2$, $X = 2$: rechterlid $\frac{4\cdot4 -
2\cdot1\cdot3}{9} = \frac{10}9$; linkerlid $\frac1{(2-1)^2} +
\frac1{(2+1)^2} = 1 + \frac19 = \frac{10}9$.

\textbf{21.} $\frac1{\binom nk} = \frac{k!\,(n-k)!}{n!} =
\frac{k!}{(n-k+1)(n-k+2)\cdots n}$, een product van $k$ opeenvolgende
gehele getallen in de noemer. Substitueren van $j = n - k + 1$ (zodat
$j$ over $\N^*$ loopt als $n$ vanaf $k$ loopt):
\[
\sum_{n=k}^{\infty}\frac1{\binom nk}
= k!\sum_{j=1}^{\infty}\frac1{j(j+1)\cdots(j+k-1)}
= \frac{k!}{(k-1)\,(k-1)!} = \frac{k}{k-1} ,
\]
door vraag 4 toe te passen met $k - 1$ in plaats van $k$ (geldig
aangezien $k - 1 \geq 1$).

\textbf{22.} Voor $k = 3$: $\frac1{\binom n3} =
\frac6{(n-2)(n-1)n}$, dus door de deelsom van vraag 3 (verschoven),
\[
\sum_{n=3}^{N}\frac1{\binom n3}
= 6\sum_{j=1}^{N-2}\frac1{j(j+1)(j+2)}
= 6\Bigl(\frac14 - \frac1{2(N-1)N}\Bigr)
= \frac32 - \frac3{(N-1)N}
\longrightarrow \frac32 ,
\]
de waarde $\frac k{k-1} = \frac32$ van vraag 21.

\textbf{23.} $n = 6$: paren $k = 1$ ($\theta_1 = \frac\pi3$,
$2\cos\theta_1 = 1$) en $k = 2$ ($\theta_2 = \frac{2\pi}3$,
$2\cos\theta_2 = -1$), plus de reële \omterm{def:b1:fractions:field}{polen} $\pm1$:
\[
\frac1{X^6-1} = \frac16\Bigl(\frac1{X-1} - \frac1{X+1}
+ \frac{X - 2}{X^2 - X + 1}
+ \frac{-X - 2}{X^2 + X + 1}\Bigr) .
\]
In $X = 0$: linkerlid $-1$; rechterlid $\frac16(-1 - 1 - 2 - 2)
= -1$: correct.

\textbf{24.} (i) De uniciteit legitimeert elke identificatie van
coëfficiënten --- afdekken, de toegevoegd complexe paargroepering van
vraag 9, en de differentiatietrucs (vragen 13, 20) berusten er alle
op. (ii) De \omterm{def:b1:complex:unity}{eenheidswortels} leverden de \omterm{def:b1:fractions:field}{polen} van $X^n -
1$, hun symmetrieën ($k \leftrightarrow n-k$), en de
halvehoekalgebra van vraag 12
(\cref{met:b1:complex:trig}). (iii) De logaritmische afgeleide
$\frac{P'}P = \sum\frac{m_i}{X-a_i}$ zette informatie over de wortels
van $P = 1 + X + \dots + X^{n-1}$ om in de numerieke sommen van de
vragen 11 en 13 --- het scharnier tussen de Delen II en III.

\textbf{25.} De splitsing is een zuiver algebraïsche identiteit, waar
voor elke waarde van de variabele tegelijk; dat is wat haar tot een
motor maakt. Gehele getallen substitueren en sommeren maakte er
telescopen van (Deel I); \omterm{def:b1:complex:unity}{eenheidswortels} substitueren en toegevoegd
complexe \omterm{def:b1:fractions:field}{polen} groeperen maakte er goniometrische identiteiten van
(Delen II en III); en pas bij de allerlaatste stap deed analyse haar
intrede --- een inknelling tussen twee gesloten vormen --- om
$\frac{\pi^2}6$ te leveren, een \omterm{def:b1:logic:statement}{uitspraak} die geen enkele eindige
substitutie kan bereiken. Deze arbeidsverdeling (algebra produceert
exacte eindige identiteiten, analyse gaat over tot de limiet) is het
sjabloon voor \cref{ch:b1:series}, waar telescoperen en vergelijken
systematisch worden, en voor \cref{ch:b1:integration}, waar elke
bouwsteen een primitieve krijgt en dezelfde splitsingen integralen
berekenen in plaats van sommen.
\end{solution}
